\documentclass[preview]{standalone}

\usepackage{amsmath}
\usepackage{amssymb}
\usepackage{bettelini}
\usepackage{stellar}
\usepackage{definitions}

\begin{document}

\id{measuretheory-egorov-lusin}
\genpage

\section{Egorov's Theorem}

\begin{snippettheorem}{egorov-theorem}{Egorov's theorem}
    Let \((X,\mathcal M,\mu)\) be a finite measure space, and let
    \(f_n,f\) be measurable functions such that \(f_n\to f\) almost
    everywhere. For every \(\varepsilon>0\), there is
    \(E\in\mathcal M\) such that
    \[
        \mu(E)<\varepsilon
    \]
    and \(f_n\to f\) uniformly on \(X\setminus E\).
\end{snippettheorem}

\begin{snippet}{egorov-theorem-interpretation}
    On a finite measure space, almost-everywhere convergence becomes uniform
    after discarding a set of arbitrarily small measure.
\end{snippet}

\begin{snippetexample}{egorov-theorem-infinite-measure-counterexample}{Egorov's theorem can fail on an infinite measure space}
    On \(\realnumbers\), let
    \[
        f_n=\chi_{[-n,n]}.
    \]
    Then \(f_n\to1\) pointwise. If \(E\) has finite measure, its complement
    is unbounded. For every \(n\), there is \(x\notin E\cup[-n,n]\), and
    therefore
    \[
        |f_n(x)-1|=1.
    \]
    The convergence cannot be uniform on \(\realnumbers\setminus E\).
\end{snippetexample}

\section{Lusin's Theorem}

\begin{snippettheorem}{lusin-theorem}{Lusin's theorem}
    Let \(f\colon\realnumbers\fromto\realnumbers\) be measurable and finite
    almost everywhere. For every measurable set
    \(A\subseteq\realnumbers\) of finite measure and every
    \(\varepsilon>0\), there is a compact set \(K\subseteq A\) such that
    \[
        \mu(A\setminus K)<\varepsilon
    \]
    and \(\restr{f}{K}\) is continuous. Equivalently, one can find a
    continuous function \(g\colon\realnumbers\fromto\realnumbers\) agreeing
    with \(f\) on \(K\), with
    \[
        \sup_{\realnumbers}|g|
        \leq\operatorname*{ess\,sup}_{A}|f|.
    \]
\end{snippettheorem}

\begin{snippetcorollary}{continuous-functions-dense-in-l-one-compact-interval}{Continuous functions are dense in \(L^1([0,1])\)}
    The space \(\mathcal C^0([0,1])\), equipped with the integral norm, is
    dense in \(L^1([0,1])\).
\end{snippetcorollary}

\begin{snippetproof}{continuous-functions-dense-in-l-one-compact-interval-proof}{continuous-functions-dense-in-l-one-compact-interval}{Continuous functions are dense in \(L^1([0,1])\)}
    First approximate an arbitrary \(L^1\) function by a bounded measurable
    function, truncating its values. Lusin's theorem then gives a compact set
    \(K\subseteq[0,1]\), with complement of arbitrarily small measure, on
    which the truncated function agrees with a continuous function \(g\).
    The truncation bounds both functions, so the integral of their difference
    on \([0,1]\setminus K\) is arbitrarily small. Combining the two
    approximations proves density.
\end{snippetproof}

\begin{snippet}{compactly-supported-continuous-functions-dense-on-real-line}
    On \(\realnumbers\), the analogous dense class is
    \(\mathcal C_c(\realnumbers)\), the continuous functions with compact
    support. One first truncates the tails of an \(L^1\) function and then
    applies the finite-measure approximation on a bounded interval.
\end{snippet}

\end{document}
